\section{Introduction}\label{sec:intro}

\subsection{The family, and the question this manuscript answers}\label{sec:intro-family}

Weil's criterion turns the Riemann hypothesis into a positivity statement: a
certain quadratic form $Q_L$, built from the primes below $e^{L}$ and an
archimedean term, is positive semidefinite for every window length $L$ if and
only if every nontrivial zero of $\zeta$ lies on the critical line. The
\emph{shifted} family \cite{WilsonLines} replaces that single form by a
one-parameter deformation. For $\omega>0$ put $s=\frac12+p$ and
\[
 K_\omega(p)=\frac{\xi(s-\omega)}{\xi(s+\omega)} .
\]
The functional equation makes $\lvert K_\omega\rvert=1$ on the imaginary axis, so
$K_\omega$ is an all-pass filter; its causal realization $V_\omega$ compressed to
a window of length $L$ is a contraction for every $L$ precisely when no zero lies
further than $\omega$ from the critical line; and the defect of that compression
is, to first order in $\omega$, twice the Weil form. The Riemann hypothesis is
therefore the assertion that the criterion holds for every $\omega>0$, and the
family interpolates between something and it.

The parent manuscript \cite{Markov} factored the transfer. Writing
$\xi(u)=\frac12u(u-1)\Lambda(u)$, the ratio
$\Kt_\omega=\Lambda(s-\omega)/\Lambda(s+\omega)$ is completely monotone --- the
Laplace transform of an explicit positive measure, a Beta-distributed continuous
delay convolved with a comb at the integers --- and what is left is a rational
all-pass factor forced by unimodularity. It also found that every piece of
$\Kt_\omega$ is $d$-dimensional as a formula, for $d=2\omega$: the archimedean
factor is a Riesz integral over $\R^{d}$, the comb's weights are the Jordan
totient $J_d$, and the residue at the Blaschke pole is the volume term of a
lattice of rank $d+1$ \cite{HalfDimension,OpeningNote,WhichDimension}. The
question that opened the present investigation is the obvious one.

\begin{problem}\label{prob:main}
Is $d=2\omega$ a dimension, or a parameter?
\end{problem}

This manuscript answers it: \textbf{a parameter}. The route runs through four
statements, and the fourth is the one worth keeping.

\subsection{What is proved}\label{sec:intro-results}

\paragraph{The transfer is an invariant of the rank
(Section~\ref{sec:rank}).} For a covolume-one lattice $L\subset\R^{m}$ let
$\Est_L(s)=\sum_{v\ \mathrm{prim}}\lvert v\rvert^{-2s}$ be the primitive Epstein
zeta. Poisson summation and the factorization $Z_L=\zeta(2s)\Est_L$ give
$\Lambda(2s)\Est_L(s)=\Lambda(m-2s)\Est_{L^{*}}(m/2-s)$, so the reflection ratio
is independent of $L$ and depends on the rank alone --- and with $m=d+1$ and
$2s=p+b$ that ratio is exactly $\Kt_\omega$. So the transfer is the scattering
matrix of the unimodular lattices of rank $d+1$, at \emph{every} integer $d\ge1$,
and its own reflection $p\mapsto-p$ is the Epstein functional equation. The
family's range $\omega\in(0,\frac12]$ is the first gap $m\in(1,2]$ of that
sequence, with the modular surface at its right endpoint.

\paragraph{The wall is the sign of $a$ (Section~\ref{sec:wall}).} With
$a=\frac12-\omega=\frac{2-m}2$ and $c=-a$, three things happen at $a=0$, that is
at $m=2$. The second pole of $\Kt_\omega$ crosses into the right half-plane ---
and it is the \emph{archimedean} factor's pole, the one from $\Gamma$, while the
first is the comb's, the one from $\zeta$ (Table~\ref{tab:poles}). The rational
factor $R_\omega=B_b\cdot\frac{p+a}{p-a}$ turns from a ratio $B_b/B_a$ with a
right half-plane pole into the inner product $B_bB_c$
(Proposition~\ref{prop:twoblaschke}). And complete monotonicity \emph{survives}:
$\Kh_\omega=\frac{p+a}{p-a}\Kt_\omega$ is completely monotone at every real
$\omega>0$, by two groupings of the Gamma factors that are exchanged at the wall,
the second being
$\Kh^\Gamma_\omega=\pi^\omega\frac{\Gamma(z+c)}{\Gamma(z+\omega)}\frac{\Gamma(z+1)}{\Gamma(z+c+1)}$
with $z=\frac{p+a}2$ (Proposition~\ref{prop:cmuncapped}). The decomposition
continues past the family's range with its positivity intact.

\paragraph{But the criterion does not (Section~\ref{sec:criterion}).} The poles
of $K_\omega$ sit at $p=\rho-b$, so they leave the closed right half-plane as
soon as $a<0$. For $\omega\ge\frac12$, therefore, every compression
$V_{\omega,L}$ is a contraction \emph{unconditionally}, on the strength of the
Euler product alone (Proposition~\ref{prop:vacuity}): the criterion is vacuous
there, and has arithmetic content on exactly $\omega\in(0,\frac12)$, where it is
equivalent to a zero-free strip of half-width $\omega$
(Corollary~\ref{cor:content}). So the Euler product is the criterion's
\emph{value} at rank two, and the Riemann hypothesis is its \emph{derivative} at
rank one, where $\Est_{\Z}\equiv2$, $\Kt_0=1$, $V_0=I$, and the first-order
defect is the Weil form (Table~\ref{tab:endpoints}). The family runs from the
easiest theorem about $\zeta$ to the hardest conjecture, and it runs downward in
rank.

\paragraph{Inside the gap there is no lattice, and the reason is Lagrange
(Section~\ref{sec:pointcount}).} The point count of $\Z^{m}$ has a unique
interpolation, the $q$-coefficients of $\theta^{m}$, polynomial in $m$. Writing
$\theta^{m}=\sum_j\binom mj(\theta-1)^{j}$ and letting $R_j(n)$ count ordered
representations of $n$ as $j$ \emph{positive} squares, one has
$r_m(n)=\sum_j\binom mj2^{j}R_j(n)$ with $R_j(n)=0$ for $j>n$. On the gap
$(k,k+1)$ the binomial is positive through $j=k+1$ and first negative at
$j=k+2$ (Proposition~\ref{prop:signlaw}), so $r_m(n)>0$ for free when $n\le k+1$
(Proposition~\ref{prop:free}) and the first coefficient that can be \emph{forced}
negative sits at $N_k$, the least integer needing $k+2$ positive squares. Those
are $N_0=2$, $N_1=3$, $N_2=7$ --- \emph{and nothing at all for $k\ge3$, because
Lagrange's four-square theorem bounds that least number by four}
(Propositions~\ref{prop:Nk} and~\ref{prop:forced}). In particular
$r_m(3)=\frac43m(m-1)(m-2)<0$ throughout $m\in(1,2)$: at fractional $d$ there is
no positive measure and hence no point set.

\paragraph{And the two failures are unrelated (Section~\ref{sec:position}).}
Both the criterion's content and the point count's failure live on $m\in(1,2)$,
which invites the guess that they are one fact. They are not. The point count
fails on $(2,3)$ as well, where the criterion is already vacuous; and at $m=3$
the point count is completely monotone while the transfer's rational factor is
not (Corollary~\ref{cor:noequiv}). The agreement is the first gap and nothing
more: the criterion's window has its endpoints at the two smallest ranks, and the
first gap of the rank sequence is $(1,2)$ by definition
(Remark~\ref{rem:forcedagreement}). What separates the two sides is where the
rank sits. In the transfer it is always an \emph{exponent} ---
$(2\sinh\tau)^{\omega-1}$ in the archimedean kernel,
$(1-e^{-t})^{\beta-\alpha-1}$ in each Beta density, $q^{-d}$ in each Euler factor
of the comb --- so each positivity is a \emph{local} inequality, at one delay or
at one prime, which continuation in the parameter cannot break. In the point
count it is a \emph{binomial index}, where positivity is a global cancellation
(Proposition~\ref{prop:position}). The transfer's positivity was therefore never
inherited from a lattice, and the lattice's failure leaves no trace in it; in
particular there is no factorization of the transfer exhibiting which factor
loses positivity at fractional $d$, because none does
(Corollary~\ref{cor:nofactorization}).

\subsection{What this is not}\label{sec:intro-not}

Nothing here constructs a source, proves any positivity of the Weil form, or
assumes the Riemann hypothesis; no statement below is conditional on it. Nothing
here is a new bound on the zeros: Proposition~\ref{prop:vacuity} is a statement
that a criterion is empty, not that anything is true of $\zeta$ beyond the Euler
product. The contraction numbers of Section~\ref{sec:calibration} are a
calibration of an inherited instrument at shifts where the answer is known in
advance, and they are one-sided in the direction that does \emph{not} certify
contraction (Remark~\ref{rem:direction}). The group-theoretic language of
Corollary~\ref{cor:oneobject} is the standard reading of verified formulas and is
labelled as such. The two statements about sums of squares in
Remark~\ref{rem:lagrange} are conjectures and are used nowhere.

Section~\ref{sec:conventions} fixes notation --- including a warning,
Remark~\ref{rem:xiwarning}, about two incompatible meanings of $\xi$ in this
program's own documents --- and lists what is inherited from \cite{Markov}.
Appendix~\ref{app:numerics} describes the six check programmes and what each
asserts; Appendix~\ref{app:provenance} records how this manuscript was drafted.
